\chapter{Modul Sumber Daya Manusia (HR)}

Modul \textit{Human Resources} (HR) bertugas mencatat dan mengelola data seluruh karyawan perusahaan, dari informasi pribadi, riwayat jabatan, hingga rekening bank untuk penggajian.

\section{Konsep Dasar \& Matriks Peran}
Data karyawan di sistem ini sangat rahasia (\textit{confidential}). Gaji dan informasi pribadi tidak boleh bisa dilihat sembarang orang.

\textbf{Pihak yang Terlibat (Role Matrix):}
\begin{itemize}
    \item \textbf{HR Staff / HR Manager}: Satu-satunya peran yang berhak menambah, mengedit, dan melihat daftar lengkap karyawan beserta detail gaji.
    \item \textbf{System Administrator}: Dapat mereset kata sandi akun karyawan jika mereka lupa.
    \item \textbf{Karyawan Umum}: (Jika fitur \textit{Employee Self-Service} diaktifkan nantinya) hanya dapat melihat data dirinya sendiri.
\end{itemize}

\section{Skenario Dunia Nyata \& Alur Kerja}
Misalkan ada karyawan baru bernama Budi yang masuk bekerja hari ini.
\begin{enumerate}
    \item \textbf{Fase 1 (Pendaftaran):} Staf HR memasukkan data diri Budi, menentukan jabatannya di departemen Produksi, serta memasukkan nomor rekening Budi.
    \item \textbf{Fase 2 (Absensi \& Cuti):} Seiring waktu, Budi bekerja, absen, atau mengambil cuti (tercatat di relasi \textit{Attendance} dan \textit{Leave}).
    \item \textbf{Fase 3 (Penggajian):} Di akhir bulan, HR mencetak Slip Gaji Budi (PDF) dan memberikannya.
    \item \textbf{Fase 4 (Reset Password):} Suatu hari Budi lupa kata sandi aplikasinya. HR mengeklik tombol "Reset Password Akun" untuk Budi.
\end{enumerate}

\section{Panduan Penggunaan (Step-by-Step)}

\subsection{1. Mendaftarkan Karyawan Baru}
\begin{enumerate}
    \item Buka menu \textbf{Employees} (Karyawan).
    \item Klik \textbf{New Employee}.
    \item Pada bagian \textbf{Data Pribadi}, isi selengkap mungkin:
    \begin{itemize}
        \item \textbf{NIK:} Nomor Induk Karyawan (wajib unik).
        \item \textbf{Nama Lengkap}, \textbf{Email}, \textbf{Nomor Telepon}, \textbf{Tanggal Lahir}, dan \textbf{Jenis Kelamin}.
        \item \textbf{Foto:} Anda bisa mengunggah file gambar (\textit{upload}) maksimal 2MB.
        \item \textbf{Alamat} domisili.
    \end{itemize}
    \item Pada bagian \textbf{Data Kepegawaian}, tentukan:
    \begin{itemize}
        \item \textbf{Departemen} dan \textbf{Jabatan}: (Catatan: Daftar jabatan hanya akan muncul jika departemen sudah dipilih).
        \item \textbf{Tipe Karyawan:} \textit{Permanent} (Tetap), \textit{Contract} (Kontrak), \textit{Probation}, atau \textit{Intern} (Magang).
        \item \textbf{Tanggal Masuk} dan \textbf{Tanggal Berakhir} (khusus untuk karyawan tidak tetap).
        \item \textbf{Gaji Pokok:} Gaji *basic* dalam Rupiah.
    \end{itemize}
    \item Lengkapi \textbf{Data Bank} (Nama Bank, Nomor Rekening, Nama Pemilik) untuk keperluan transfer gaji.
    \item Lengkapi \textbf{Kontak Darurat} (Nama dan Nomor Telepon).
    \item Klik \textbf{Create}. Data Budi resmi tersimpan.
\end{enumerate}

\begin{figure}[H]
    \centering
    \includegraphics[width=0.8\textwidth]{placeholder_form_employee.png}
    \caption{Formulir Pendaftaran Karyawan Baru}
    \label{figure:6.form_employee}
\end{figure}

\subsection{2. Mencetak Slip Gaji (Export PDF Slip)}
\begin{enumerate}
    \item Buka tabel daftar \textbf{Employees}.
    \item Cari nama karyawan yang diinginkan.
    \item Di kolom Action paling kanan, temukan tombol berikon Dokumen (\textbf{Export PDF Slip}).
    \item Sistem akan secara otomatis mengunduh file PDF berisikan slip gaji terakhir dari karyawan tersebut.
\end{enumerate}

\subsection{3. Mereset Kata Sandi Karyawan}
Terkadang karyawan lupa *password* untuk masuk ke dalam sistem ERP. HR atau Admin dapat membantunya.
\begin{enumerate}
    \item Buka daftar \textbf{Employees}.
    \item Cari nama karyawan tersebut.
    \item Klik tombol berikon Kunci Kuning (\textbf{Reset Password Akun}). Tombol ini hanya muncul jika karyawan tersebut sudah punya akun \textit{User} yang tertaut.
    \item Masukkan \textbf{Password Baru} dan ketik ulang di \textbf{Konfirmasi Password}.
    \item Klik \textit{Confirm}. Saat karyawan tersebut masuk menggunakan kata sandi baru ini, sistem akan memaksanya (\textit{force change}) untuk mengganti lagi *password* tersebut demi keamanan.
\end{enumerate}

\section{Penanganan Masalah (Edge Cases)}
\begin{itemize}
    \item \textbf{Error saat mencetak Slip Gaji ("Belum ada data payroll"):} Anda tidak bisa mencetak slip gaji jika belum ada data *Payroll* (Penggajian) bulan ini untuk karyawan tersebut yang digenerate oleh sistem di menu *Payroll*.
    \item \textbf{Jabatan tidak bisa dipilih:} Pastikan Anda sudah memilih \textbf{Departemen} di *dropdown* atasnya. Jabatan bergantung/terkunci pada Departemen.
\end{itemize}
